\section{The critical mass of paths in a box}\label{sec:full-box}
The later proofs extract a special subpath and sum freely over the two
outside pieces. We therefore need a bound that sums \emph{all} lengths
and endpoints in the containing box, with one fixed polynomial exponent.
The one-arc estimate in Theorem~\ref{09_fast_paths:one-arc-bound} supplies
that bound after a geometric reduction. Its uniformity when one group of
bands stays fixed while the other grows is used in the last step.

\begin{proposition}[Full-box susceptibility]\label{09_fast_paths:full-box}
Let \(S(H)\) be the maximum, over starting vertices in a box of side
\(H\), of the critical mass of walks confined to that box, with
arbitrary endpoint and length. There are constants \(C_0,C_1\) such
that \(S(H)\le C_0(1+H)^{C_1}\).
\end{proposition}
\begin{proof}
\emph{The suffix test.}
We first choose a small constant \(a>0\), and later choose a sufficiently
large \(W\), once and independently of \(H\). Put
\(m=\lfloor aW\rfloor\). Walks shorter than \(m\) cost
\(\exp(O(W))\). Consider the finitely many triangular bases, including
reflections, whose upward side directions span a \(60^\circ\) cone.
In each basis require a westward step among the last \(m\) steps.
If this fails in one basis, there is at most one eastward and one
bisector-parallel choice at each vertex. Two consecutive parallel choices
would immediately backtrack. Thus the number of length-\(l\)
continuations avoiding west is at most \(C\varphi^l\), where
\(\varphi=(1+\sqrt5)/2\). Since \(\rho\varphi<1\), the total mass
failing at least one of these suffix tests is at most
\begin{equation}\label{eq:fullbox-suffix}
 \epsilon_W S(H),\qquad
 \epsilon_W=C(\rho\varphi)^m.
\end{equation}
The prefix ending where the suffix starts is included in \(S(H)\);
ignoring its avoidance constraints only enlarges this bound.

\smallskip\noindent\emph{Finding a cylinder arc.}
We seek a middle arc from a suitable earlier step to a westward step in
the suffix. The criterion below makes this arc eligible for the one-arc
theorem. Choosing the first such step leaves a prefix with no earlier
qualifying step; after treating the arc, we will bound the mass of these
prefixes by a constant depending only on \(W\).
Fix the endpoint \(z\), costing \(O((1+H)^2)\) choices, and scan the
step midpoints before this suffix. Let \(\ell_1,\ell_2,\ell_3\) be the
three lattice-side lines through \(z\), and put
\(d(p)=\min_i\operatorname{dist}(p,\ell_i)\). If \(d(p)>W\), the
vector \(q=z-p\) lies in one open positive cone. Write
\[
 q=u_1e_1+u_2e_2,
 \qquad u_1,u_2\ge 2W/\sqrt3,
 \qquad t=\frac{u_1}{u_1+u_2},
\]
where \(e_1,e_2\) are its adjacent unit triangular side vectors. The
coefficient bounds follow by taking distances to the two cone-side lines.
Assignment \(A\) makes \(e_1\) the upper-angle direction and is eligible
when \(t\ge r\), where \(r=1/2-\delta_1/2\). The reflected assignment
\(B\) makes \(e_2\) upper and is eligible when \(t\le1-r\). Reflection
fixes the cone bisector and reverses the transverse coordinate. In
particular, a step with negative transverse increment is westward in
\(A\), a step with positive increment is westward in \(B\), and a
bisector-parallel step is westward in neither. Call a step successful
if it is westward in some eligible assignment.

Suppose there is a success, and use the first one, with midpoint \(p\).
Choose a westward suffix step in its eligible basis, with midpoint \(p'\).
The intervening arc departs west from \(p\) and arrives east at \(p'\).
Moreover \(p'-p=q+O(m)\). The two marked triangular sides have types
\(\tau,\tau'\in\{1,2\}\); write
\(p=v+e_\tau/2\) and \(p'=v'+e_{\tau'}/2\), where \(v,v'\) are
triangular lattice vertices. The displacement between the two marked
half-sides is
\begin{equation}\label{eq:fullbox-integer-row}
 n_1e_1+n_2e_2
   =(p'-p)-\frac{e_\tau+e_{\tau'}}2
   =v'-v-e_\tau.
\end{equation}
Thus \(n_1,n_2\) are integers, since adjacent side vectors form an
integer lattice basis. They satisfy
\(n_i\ge cW-Cm-1\), and \(s=n_1+n_2\le C(H+1)\).
The change in the upper coefficient fraction is at most
\(C(m+1)/W\). Choose \(a\) small relative to \(\delta_1\), and then
\(W\) large. Both coefficients are positive, \(s\ge cW\), and the
ordinary upper fraction is at least \(1/2-\delta_1\).

Figure~\ref{fig:fullbox-cone-row} shows the cone and the marked
half-sides in this displacement calculation.
\begin{figure}[htbp]
\centering
\begin{tikzpicture}[x=1cm,y=1cm,font=\small,>=Stealth,
  row/.style={line width=1pt,line join=round},
  marker/.style={line width=1.6pt,blue!65!black}]
 \begin{scope}
  \coordinate (p) at (0,0);
  \coordinate (q) at (3,1.732);
  \fill[gray!7] (p)--(4.20,0)--(2.10,3.637)--cycle;
  \draw[->] (p)--(4.20,0) node[below] {$e_1$};
  \draw[->] (p)--(2.10,3.637) node[above] {$e_2$};
  \draw[densely dashed,gray] (2,0)--(q)--(1,1.732);
  \draw[->,line width=.8pt] (p)--(q);
  \fill (p) circle (1.4pt) node[below left] {$p$};
  \fill (q) circle (1.4pt) node[above right] {$z$};
  \draw (.48,0) arc (0:60:.48);
  \node at (.68,.78) {$60^\circ$};
  \node[rotate=30,fill=gray!7,inner sep=1.5pt] at (1.82,1.42) {$q=z-p$};
  \node[anchor=west,font=\footnotesize] at (0,4.12) {Positive lattice cone};
  \node[anchor=north] at (1.65,-.53) {$q=u_1e_1+u_2e_2$};
 \end{scope}
 \begin{scope}[xshift=6.5cm]
  \coordinate (v) at (0,0);
  \coordinate (p) at (.25,.433);
  \coordinate (w) at (.5,.866);
  \coordinate (vp) at (4.5,2.598);
  \coordinate (pp) at (5,2.598);
  \coordinate (end) at (5.5,2.598);
  \draw[marker,opacity=.28] (v)--(p);
  \draw[marker] (p)--(w);
  \draw[row] (w)--(3.5,.866)--(vp);
  \draw[marker] (vp)--(pp);
  \draw[marker,opacity=.28] (pp)--(end);
  \foreach \x/\y in {0/0,.5/.866,1.5/.866,2.5/.866,3.5/.866,4/1.732,4.5/2.598,5.5/2.598}
    \fill (\x,\y) circle (1pt);
  \fill[blue!65!black] (p) circle (1.9pt);
  \fill[blue!65!black] (pp) circle (1.9pt);
  \node[anchor=east] at (.07,.44) {$p$};
  \node[anchor=south] at (5,2.66) {$p'$};
  \node[anchor=north east] at (-.06,-.04) {$v$};
  \node[anchor=south east] at (4.44,2.67) {$v'$};
  \node[anchor=south] at (1.95,.98) {$n_1e_1$};
  \node[anchor=east] at (3.72,1.80) {$n_2e_2$};
  \node[anchor=west,text=blue!65!black] at (.59,.32) {$\tfrac12e_\tau$};
  \node[text=blue!65!black] at (5.0,2.12) {$\tfrac12e_{\tau'}$};
  \draw[->,gray!75] (5.57,2.598)--(6.10,2.598);
  \node[anchor=west,font=\footnotesize] at (6.17,2.598) {$Y$};
  \node[anchor=west,font=\footnotesize] at (0,4.12) {Marked half-sides and ordinary $X$-bands};
  \node[anchor=north,align=center] at (2.65,-.53)
    {$p'-p=\tfrac12e_\tau+n_1e_1+n_2e_2+\tfrac12e_{\tau'}$};
 \end{scope}
\end{tikzpicture}
\caption{The row construction in a positive $60^\circ$ lattice cone.
On the right, an example canonical order uses $n_1=3$ ordinary sides in
direction $e_1$, followed by $n_2=2$ in direction $e_2$; the marked
sides have types $\tau=2$ and $\tau'=1$. The segment from $p$ to
$p'$ contains one half of each marked side, so subtracting these
halves leaves the integer displacement $n_1e_1+n_2e_2$.
The faint halves complete the marked sides but do not belong to this
segment. This is a row cut, not the extracted honeycomb arc; the arc
may cross ordinary parts of the row. The angles drawn are Euclidean
lattice angles, not the spectral parameters.}
\label{fig:fullbox-cone-row}
\end{figure}


Use these integers as the ordinary \(X\)-band counts in the row
\([p,X,p',Y]\), in a fixed canonical order. We do not sum over band
orders. Choose all \(b\) ordinary \(Y\)-bands lower, with
\(b\ge K(H+s+1)\), adjusting by at most one to make \(s+b+2\) even.
We may take \(b\le C(H+s+1)\). The period vector, including the two
full marked sides, has projection at least \(cb\) on the cone bisector.
Taking \(K\) large makes its norm larger than the box diameter with the
bounded endpoint margin. Projection to this cylinder is consequently
injective on the box. The selected arc has no repeated marked port, and
its unused marked half-edges are vacant. Ordinary crossings of the row
are permitted. The physical row therefore gives exactly an arc in
Theorem~\ref{09_fast_paths:one-arc-bound}, with unchanged vertex weight.

The anchored port, basis, two marked-side types and two ordinary band
counts determine the row and its lift once \(b\) has been chosen;
there are only polynomially many such data in the box. Choose \(W\)
also so that \(s\ge cW\) implies \(L_X\ge L_*\). The arc mass is then
at most \(b^C\le C(1+H)^C\), while its suffix remnant costs
\(\exp(O(W))\). Cutting at the two midpoints changes the vertex-path
weight by bounded factors. It remains to bound the mass of the
unsuccessful prefix uniformly in \(H\) and \(z\). If no step succeeds,
the same prefix estimate applies to the whole part before the suffix.

\smallskip\noindent\emph{Summing the unsuccessful prefixes.}
We first record a quantitative strip estimate. From one fixed start,
the total free-end mass of all lengths in a straight strip of width
\(O(W)\), parallel to a lattice side, is at most
\begin{equation}\label{eq:fullbox-strip}
 U(W)=C\exp(CW^{3/4}).
\end{equation}
Indeed, split at the first global minimum in the normal direction,
reverse the first tail, and retain the second. Decompose each tail
chronologically at successive last alternating extrema of its remaining
piece. The spans satisfy \(a_1\ge a_2>a_3>\cdots\) on the
one-third-band grid; each positive span occurs at most twice per tail.
By Lemma~\ref{lem:port-extension} and Theorem~\ref{thm:01-strip}, a weak
bridge of span \(j/3\) has free-end mass at most
\(f(j)=C'(1+j)^{-1/4}\). Zero-span pieces are trivial. Bounded endpoint
and cut records can be incorporated in \(C'\). Summing the span lists
for the two tails gives
\[
 \prod_{1\le j\le CW}(1+f(j))^4\le\exp(C''W^{3/4}).
\]
Anchor the minimum at one of the two honeycomb translation-class
representatives. The two complete tails determine the split time, and
the endpoint of the reversed first tail determines the unique translation
returning to the prescribed original start. Thus there is no additional
sum over lengths, split times or minimum positions. This proves
\eqref{eq:fullbox-strip}.

Next consider a maximal run of unsuccessful midpoints with \(d(p)>W\).
Its cone is fixed. The cone coefficients change by a bounded amount per
step and remain at least \(cW\), so consecutive values satisfy
\( |t_{i+1}-t_i|\le C/W\). Choose \(W\) so large that this is less
than \(\delta_1/4\). A run visiting both the sole-\(A\) region
\(t>1-r\) and the sole-\(B\) region \(t<r\) would have at least two
consecutive midpoints in the overlap \([r,1-r]\). There both assignments
are eligible. An unsuccessful step then has neither positive nor negative
transverse increment and must be parallel to the cone bisector. Every
honeycomb vertex has exactly one incident edge in that direction, so two
consecutive such steps backtrack. This is impossible. Hence the run
visits at most one sole-choice region; its corresponding assignment
(or either assignment if there is no sole-choice region) is eligible
throughout. Every step avoids west in this one fixed assignment. Its
length-\(l\) mass is therefore at most \(C(\rho\varphi)^l\), by the
same Fibonacci count used in \eqref{eq:fullbox-suffix}.

Retain only those maximal runs with \(\max d(p)\ge2W\). Since \(d\)
changes by a bounded amount per step, every retained interior run has
length at least \(cW\). Only the first and last retained runs can lack
this lower bound. Summing the lengths of each interior run costs at most
\(C e^{-c'W}\); either end run has bounded total mass. All other
midpoints lie within \(2W\) of one of the three lines. Outside a disk
of radius \(C_0W\) about \(z\), these three strips are disjoint and
separated by more than a step. A remaining piece stays in one straight
strip until it returns to the disk or begins a retained run. The disk
contains \(O(W^2)\) vertices, so self-avoidance permits only \(O(W^2)\)
disk steps, which we enumerate individually.

With \(j\) retained interior runs there are at most
\(j+C_1W^2+C_1\) strip pieces. Record their ordered types, strip labels
and fixed assignments, each with boundedly many choices. Intermediate
positions and lengths are already summed by enumerating each piece from
its current start; there is no factor depending on \(H\) at a join.
By \eqref{eq:fullbox-strip}, the unsuccessful-prefix mass is bounded by
\begin{equation}\label{eq:fullbox-prefix}
 P_W\le N(W)e^{CW^2}
       \sum_{j\ge0}\bigl[K U(W)e^{-c'W}\bigr]^j,
 \qquad N(W)=C U(W)^{C_1W^2+C_1}.
\end{equation}
The factor \(N(W)\) also absorbs the at most two end runs. Its possible
size \(\exp(O(W^{11/4}))\) is harmless: \(W\) is fixed. The series
converges for sufficiently large \(W\), since \(W^{3/4}=o(W)\).

After choosing \(a\), choose \(W\) once to satisfy all the preceding
coefficient, fraction, strip and one-arc requirements, and also
\(\epsilon_W<1/2\). Boxes with \(H\le CW\) have a finite uniformly
bounded number of self-avoiding walks and can be absorbed in a constant
depending on \(W\). For larger boxes, sum the endpoint, row and lift
data, apply the prefix bound and the one-arc estimate, and then add
\eqref{eq:fullbox-suffix}. This gives, with a fixed exponent \(K_0\),
\[
 S(H)\le A_W(1+H)^{K_0}+\epsilon_W S(H).
\]
The box has finitely many vertices, so \(S(H)<\infty\). Absorbing the
last term proves the proposition.
\end{proof}
